% Integration-ready research section. Parent supplies theorem environments,
% amsmath, amssymb, amsthm, and hyperref. Bibliography keys suggested below.
\section{A square-root transition in the late reflected residues}
\label{moderate:sec}

This section concerns a different simultaneous limit from the transition
$n/(m+1)-\log m=O(1)$ for the moments themselves. Here the late residue
index $k$ grows, and the reflected exponent is allowed to satisfy
$m=O(\sqrt{k})$. The fixed-$m$ asymptotic in the manuscript does not
cover this limit. The novelty assertion here is relative to the pinned
repository: no claim of priority over all literature is made. We prove a Gaussian damping factor, give its first
correction, and then obtain a sharp truncation theorem for a growing
reflected exponent.

Put
\[
 f(x)=\log\Gamma(x),\quad B(x)=\log\Gamma(1-x),\quad
 M_{n,m}=\int_0^1 f(x)^n f(1-x)^m\,dx,
\]
and retain the manuscript's residues
\begin{equation}\label{moderate:eq:A}
 A_{k,m}=\frac1k[x^{k-1}]\Gamma(1+x)^k B(x)^m.
\end{equation}
The parameters $k,m$ are integers with $k\ge1$, $m\ge0$.
The coefficients vanish when $k\le m$. In the meromorphic transform
$F_m(z)=\int_0^1\Gamma(x)^zB(x)^m\,dx$, the actual residue is
$\operatorname*{Res}_{z=k}F_m(z)=-kA_{k,m}$, so its sign is the
opposite of the coefficient sign. The standard Taylor series
for log gamma used below follows from DLMF, Section~5.7\cite{moderate:DLMF}.
Lagrange inversion and the underlying saddle and singularity methods are
classical\cite{moderate:Gessel,moderate:FS}; the moving-parameter estimates
are proved explicitly here.
Define
\[
 \phi(u)=\Gamma(1-u),\quad a(u)=-\log\Gamma(1+u),\quad D=u\frac d{du}.
\]
Let $\tau\in(0,1)$ be the unique solution of
$D\log\phi(\tau)=1$, and put
\begin{equation}\label{moderate:eq:constants}
 \rho=\frac{\tau}{\phi(\tau)},\quad
 \Lambda=-\log\rho,\quad b=a(\tau),\quad
 v=D^2\log\phi(\tau),\quad
 C=\frac{\tau}{\sqrt{2\pi v}}.
\end{equation}
All these constants are positive. Use a local real branch of $\log a$
at $\tau$ and set
\begin{equation}\label{moderate:eq:derivatives}
 r_j=D^j\log a(\tau),\qquad
 \kappa_j=D^j\log\phi(\tau),\qquad
 \delta_\Gamma=\frac{r_1^2}{2v}.
\end{equation}
Thus $\kappa_2=v$. The subscript on $\delta_\Gamma$ distinguishes this
constant from a small contour neighborhood.
The damping constant is strictly positive. Indeed
$D\log\phi(1/2)=\gamma/2+\log2<1$ and
$\psi(3/2)=2-\gamma-2\log2>0$, using
$\gamma<3/5$ and $\log2<7/10$. Hence $\tau>1/2$,
$\psi(1+\tau)>0$, and
$r_1=-\tau\psi(1+\tau)/b<0$.
Numerically,
\[
 b=0.1206254661326052\ldots,\quad
 r_1=-0.1684096130814827\ldots,\quad
 \delta_\Gamma=0.00624269203454726\ldots.
\]
These decimals only illustrate the constants; their defining equations
are used in the proof.

\begin{theorem}[Uniform late-residue transition]\label{moderate:thm:late}
Fix $L<\infty$. Uniformly for integers $0\le m\le L\sqrt{k}$ as
$k\to\infty$, with $\xi=m/\sqrt{k}$,
\begin{equation}\label{moderate:eq:late}
 A_{k,m}=(-1)^{k+m-1}C b^m\rho^{-k}k^{-3/2}
 e^{-\delta_\Gamma\xi^2}
 \left(1+\frac{P_1(\xi)}{\sqrt{k}}+O_L(k^{-1})\right),
\end{equation}
where
\begin{align}\label{moderate:eq:P1}
 P_1(\xi)={}&\left[-\frac{r_1+r_2/2}{v}
             +\frac{\kappa_3r_1}{2v^2}\right]\xi\notag\\
 &+\left[\frac{r_2r_1^2}{2v^2}
             -\frac{\kappa_3r_1^3}{6v^3}\right]\xi^3.
\end{align}
For every fixed $J\ge0$ there are explicitly computable real
polynomials $P_0,\ldots,P_J$, with $P_0=1$, such that the parenthesis
in \eqref{moderate:eq:late} equals
\[
 \sum_{j=0}^J k^{-j/2}P_j(\xi)
       +O_{L,J}(k^{-(J+1)/2}).
\]
In particular the coefficient signs are uniformly $(-1)^{k+m-1}$ for all
sufficiently large $k$ in this moving range. If $m/\sqrt{k}\to\xi_0$,
then the ratio of $A_{k,m}$ to its fixed-$m$ leading expression converges
to $e^{-\delta_\Gamma\xi_0^2}$.
\end{theorem}

\begin{proof}
Replacing $x$ by $-u$ in \eqref{moderate:eq:A} gives
\[
 A_{k,m}=\frac{(-1)^{k+m-1}}{k}
        [u^{k-1}]\phi(u)^k a(u)^m.
\]
The series of $\phi$ has strictly positive coefficients, because
\[
 \log\phi(u)=\gamma u+\sum_{j\ge2}\frac{\zeta(j)}j u^j.
\]
Consequently
\[
 Q(t)=\frac{\phi(\tau e^{it})}{\phi(\tau)}e^{-it}
\]
satisfies $|Q(t)|<1$ for $0<|t|\le\pi$, and
\[
 \log Q(t)=-vt^2/2-i\kappa_3t^3/6+\kappa_4t^4/24+O(t^5).
\]
Cauchy's coefficient formula is
\begin{equation}\label{moderate:eq:cauchy}
 A_{k,m}=\frac{(-1)^{k+m-1}\tau\rho^{-k}b^m}{2\pi k}
 \int_{-\pi}^{\pi}Q(t)^k e^{it}
       \left(\frac{a(\tau e^{it})}{b}\right)^m\,dt.
\end{equation}
No logarithm of $a$ on the full circle is required: $a^m$ is an analytic
integer power. Choose a small fixed $\epsilon>0$ such that $a$ is nonzero
near the arc $|t|\le\epsilon$. Outside that arc, compactness gives a
bound $O(q^k H^m)$ with $q<1$ and a fixed $H$. Since $m\le L\sqrt{k}$,
this is exponentially small in $k$.

On the small arc use the logarithm with value $\log b$ at zero.
Real Taylor coefficients imply
\[
 \log\frac{a(\tau e^{it})}{b}
 =ir_1t-r_2t^2/2-ir_3t^3/6+O(t^4).
\]
Its real part is $O(t^2)$. After reducing $\epsilon$ if necessary,
$|Q(t)|^k|a(\tau e^{it})/b|^m\le e^{-c kt^2}$ with $c>0$, uniformly
in the stated range once $k$ is large. Thus the portion
$k^{-2/5}\le|t|\le\epsilon$ is $O(e^{-c k^{1/5}})$.
Set $t=z/\sqrt{k}$ on the remaining arc. Its integrand is
\begin{align*}
 e^{-vz^2/2+ir_1\xi z}\bigg[1+\frac1{\sqrt{k}}
 \left(iz-\frac{i\kappa_3z^3}{6}-\frac{\xi r_2z^2}{2}\right)
 +O_L\left(k^{-1}(1+|z|^6)\right)\bigg].
\end{align*}
The error is bounded by an integrable polynomial times a fixed Gaussian;
this follows from the analytic Taylor remainder and the preceding
Gaussian bound. Extending the $z$ interval to $\mathbb R$ has an
exponentially small cost. The leading integral is
\[
 \int_{\mathbb R}e^{-vz^2/2+ir_1\xi z}\,dz
  =\sqrt{\frac{2\pi}{v}}e^{-r_1^2\xi^2/(2v)}.
\]
It is bounded away from zero relative to $\sqrt{2\pi/v}$ for
$0\le\xi\le L$, so all subsequent error estimates remain relative
and uniform.

Gaussian differentiation gives the normalized moments
\[
 \langle z\rangle=\frac{ir_1\xi}{v},\qquad
 \langle z^2\rangle=\frac1v-\frac{r_1^2\xi^2}{v^2},\qquad
 \langle z^3\rangle=i\left(\frac{3r_1\xi}{v^2}
                          -\frac{r_1^3\xi^3}{v^3}\right).
\]
Substituting these into the term of order $k^{-1/2}$ gives
\eqref{moderate:eq:P1}. For any higher fixed order, expand the two
analytic logarithms and $e^{it}$ to the corresponding finite order,
then integrate the resulting polynomial against this same complex
Gaussian. Gaussian domination bounds the remainder. Differentiation
of the Gaussian Fourier transform shows that every normalized
coefficient is a polynomial in $\xi$. This proves the full expansion
and all assertions of the theorem.
\end{proof}

A completely explicit coefficient prescription is available. With
$\varepsilon=k^{-1/2}$, expand
\begin{equation}\label{moderate:eq:algorithm}
 \exp\left\{i\varepsilon z+
 \sum_{j\ge3}\frac{\kappa_j(iz)^j}{j!}\varepsilon^{j-2}
 +\xi\sum_{j\ge2}\frac{r_j(iz)^j}{j!}\varepsilon^{j-1}\right\}
 =\sum_{j\ge0}T_j(z,\xi)\varepsilon^j.
\end{equation}
Then $P_j(\xi)$ is the integral of $T_j$ against
$e^{-vz^2/2+ir_1\xi z}$, divided by the leading Gaussian integral.
Only finitely many terms are needed for each $j$.
As a consistency check, the fixed-$m$ coefficient in the manuscript is
\[
 \beta_m=\beta_0+
 m\left[-\frac{r_1+r_2/2}{v}+\frac{\kappa_3r_1}{2v^2}\right]
 -\delta_\Gamma m^2.
\]
The quadratic term is exactly the first Taylor coefficient of the
Gaussian damping factor. The uniform theorem resums that part when
$m^2/k$ is no longer negligible.

\subsection{Uniform optimal truncation with a growing reflected exponent}

\begin{theorem}[A moving-exponent half-term law]\label{moderate:thm:half}
Fix $L,E<\infty$. Let $n,N\to\infty$ through positive integers,
$0\le m\le L\sqrt{N}$, and $|\eta|\le E$, where
$\eta=N\Lambda-n$. Then
\begin{align}\label{moderate:eq:half}
 \frac{M_{n,m}/n!-\sum_{k=1}^{N-1}A_{k,m}/k^n}{A_{N,m}/N^n}
 =\frac12+
 \frac{3/2-\delta_\Gamma m^2/N-\eta}{4N}
 +O_{L,E}(N^{-3/2}).
\end{align}
The denominator is nonzero for all sufficiently large $N$, uniformly
in the stated range. In particular the error is asymptotic to half the
first omitted term, and its eventual sign is $(-1)^{N+m-1}$.
The theorem supplies a proved growing range; it does not assert this
range is maximal.
\end{theorem}

\begin{proof}
We give the contour and tail details because neither a formal growing
substitution in the fixed-$m$ formula nor an infinite divergent Taylor
tail would justify this statement.
Let $h$ be the inverse germ $h=q\phi(h)$, let $g(t)=-h(-t)$, and set
\[
 W_m(x)=\int_0^x B(u)^m\,du,\qquad G_m(t)=W_m(g(t)).
\]
Lagrange inversion gives $G_m(t)=\sum_{k\ge1}A_{k,m}t^k$ near zero.
The manuscript's inverse-gamma argument shows that $g$ continues to a
slit disk of radius greater than $\rho$, whose only singularity on its
circle of convergence is at $-\rho$. For completeness, positivity of
the coefficients of $h$ gives $|h(q)|<\tau$ at every $|q|=\rho$,
$q\ne\rho$; then
$|q\phi'(h(q))|<\rho\phi'(\tau)=1$ allows analytic continuation.
At $q=\rho$ the first derivative of $u/\phi(u)$ vanishes and the second
does not, so ordinary local inversion gives a convergent square-root
expansion. Compactness produces a common slit disk. Shrink it so
$|g|<r_0<1$ there; this makes composition with every $W_m$ valid.

Choose a fixed outer radius $R>\rho$ in that disk and a small
$0<\epsilon<\Lambda$ with $\rho e^\epsilon<R$. Along its negative cut define
the real, possibly signed density
\[
 \nu_m(s)=\frac{(-1)^m}{2\pi i}
        \bigl(G_m(-s+i0)-G_m(-s-i0)\bigr),\qquad \rho<s<R.
\]
Put $\mathcal B_{N,m}=(-1)^{N+m-1}A_{N,m}$.
Cauchy contour deformation yields
\begin{equation}\label{moderate:eq:cut-coeff}
 \mathcal B_{N,m}=I_{N,m}+O(M^mR^{-N}),\qquad
 I_{N,m}=\int_\rho^R s^{-N-1}\nu_m(s)\,ds,
\end{equation}
for a fixed $M$, with an inessential enlargement of the constant.
The corresponding remainder contour has the additional kernel
$s/(s+t)$. The common factor and sign cancel after division by
$A_{N,m}t^N$. Outer-contour errors, including this division, are
exponentially small for $m\le L\sqrt{N}$ and
$\rho e^{-\epsilon}\le t\le\rho e^\epsilon$.

The jump need not be positive in this moving range. The needed absolute
bound is as follows. Write $u=(s-\rho)/\rho$. The square-root inverse
has $g(-s\pm i0)=-\tau+O(\sqrt{u})$ and a connecting path of length
$O(\sqrt{u})$. On that path
$|B(z)|\le b\exp(C_0\sqrt{u})$ for sufficiently small fixed $u$.
Since $W_m'=B^m$,
\[
 |\nu_m(\rho(1+u))|
 \le C_1 b^m\sqrt{u}\exp(C_0m\sqrt{u}).
\]
On the rest of the finite cut it is at most $C_2M^m$. Substituting
$w=Nu$, using $m\le L\sqrt{N}$ and
$(1+u)^{-N-1}\le e^{-cNu}$ on the small interval, proves for every
fixed $j\ge0$
\begin{equation}\label{moderate:eq:absolute-jump}
 \int_\rho^R \left|\frac{s-\rho}{\rho}\right|^j
    s^{-N-1}|\nu_m(s)|\,ds
 \le C_{L,j} b^m\rho^{-N}N^{-j-3/2}.
\end{equation}
The integrable majorant is a constant times
$w^{j+1/2}e^{-cw+C_0L\sqrt{w}}$.
By Theorem~\ref{moderate:thm:late}, division by $\mathcal B_{N,m}$
turns \eqref{moderate:eq:absolute-jump} into $O_{L,j}(N^{-j})$.
This absolute estimate prevents cancellation in the signed jump from
invalidating the Taylor-remainder bounds.

The first normalized cut moment is computed without expanding the
moving jump explicitly. The exact finite-cut identity gives
\[
 \int_\rho^R\frac{s-\rho}{\rho}s^{-N-1}\nu_m(s)\,ds
   =\rho^{-1}I_{N-1,m}-I_{N,m}.
\]
Use Theorem~\ref{moderate:thm:late} through order $N^{-1}$, with
relative remainder $O_L(N^{-3/2})$, both at $N-1$ and at $N$ with
the same $m$. The polynomials in that expansion are smooth for bounded
$m/\sqrt N$. Their ratio changes by $O_L(N^{-3/2})$, while the leading
factors give
\begin{align*}
 \frac{\mathcal B_{N-1,m}}{\rho\mathcal B_{N,m}}
 &=\left(\frac{N}{N-1}\right)^{3/2}
    \exp\left(-\frac{\delta_\Gamma m^2}{N(N-1)}\right)
    \bigl(1+O_L(N^{-3/2})\bigr)\\
 &=1+\frac{3/2-\delta_\Gamma m^2/N}{N}
       +O_L(N^{-3/2}).
\end{align*}
For $r=t/\rho$,
\[
 \frac{s}{s+t}=\frac1{1+r}+
            \frac{r}{(1+r)^2}u+O(u^2)
\]
uniformly on the fixed middle interval. Equations
\eqref{moderate:eq:cut-coeff}--\eqref{moderate:eq:absolute-jump} therefore give
\begin{align}\label{moderate:eq:inverse-remainder}
 \frac{G_m(t)-\sum_{k<N}A_{k,m}t^k}{A_{N,m}t^N}
 =\frac1{1+r}+
 \frac{r}{(1+r)^2}\frac{3/2-\delta_\Gamma m^2/N}{N}
 +O_L(N^{-3/2}).
\end{align}

It remains to justify averaging this local statement. The exact
layer-cake identity is
\[
 \frac{M_{n,m}}{n!}=\frac1{\Gamma(n)}
       \int_0^\infty y^{n-1}G_m(e^{-y})\,dy.
\]
After subtracting the finite sum, the ratio in
\eqref{moderate:eq:half} is the expectation of the left side of
\eqref{moderate:eq:inverse-remainder} at $t=e^{-Y}$, where $Y$ has
gamma shape $n$ and rate $N$.
Its probability outside $|Y-\Lambda|\le\epsilon$ is $O(e^{-cN})$
for bounded $\eta$, by the elementary gamma moment-generating function.
Uniform subexponential bounds on the normalized remainder suffice on
that complement. To obtain them, Cauchy's coefficient formula on
$|u|=\tau$ gives
\[
 |A_{k,m}|\le\frac{\tau}{k}\rho^{-k}H^m
\]
with a fixed $H$. Below $t=\rho e^{-\epsilon}$ sum the absolutely
convergent Taylor tail. Above $t=\rho e^\epsilon$ bound the finite
polynomial by its absolute terms. Division by the uniform lower bound
for $|A_{N,m}|$ in Theorem~\ref{moderate:thm:late} gives at most a
polynomial in $N$ times $e^{C_L\sqrt N}$ in both cases.
The remaining function term on the upper interval satisfies
$0\le G_m(t)\le m!$ for $0\le t\le1$: indeed $0\le g(t)\le1$
and $B(x)\le-\log(1-x)$, since
$\log\Gamma(2-x)\le0$ by log convexity between one and two. Thus
$W_m(g(t))\le\int_0^1[-\log(1-x)]^m\,dx=m!$.
After division its size is
at most a polynomial times
$e^{-\epsilon N}m!b^{-m}$, hence exponentially small in $N$ because
$m\log(m+1)=O_L(\sqrt N\log N)$.
The weighted outer contribution is therefore exponentially small.

Finally put $V=Y-\Lambda$. Then
\[
 \mathbb EV=-\eta/N,\quad \mathbb EV^2=O_E(N^{-1}),\quad
 \mathbb EV^3=O_E(N^{-2}),\quad \mathbb EV^4=O_E(N^{-2}).
\]
The leading kernel has expansion
$1/(1+e^{-V})=1/2+V/4-V^3/48+O(V^4)$, so its expectation is
$1/2-\eta/(4N)+O_E(N^{-2})$. The next kernel is
$e^{-V}/(1+e^{-V})^2=1/4+O(V^2)$, since its derivative vanishes
at zero. Its expectation is $1/4+O_E(N^{-1})$.
Averaging \eqref{moderate:eq:inverse-remainder} gives
\eqref{moderate:eq:half}. The signs and nonvanishing follow from
Theorem~\ref{moderate:thm:late}.
\end{proof}

\paragraph{What the theorem resolves and what it leaves open.}
The fixed-reflected-exponent warning in the manuscript is correct; it
should now be supplemented by a separate theorem for $m=O(\sqrt N)$.
The square-root scale is the first scale on which the leading late
residue changes by a nontrivial factor. It does not destroy the
half-term law, but it changes that law's first correction. The
fixed-$m$ correction $\beta_m/(4N^2)$ contains the same term
$-\delta_\Gamma m^2/(4N^2)$; the result above proves the relevant
resummation uniformly instead of substituting a growing parameter in
an uncontrolled remainder.

The next subsection proves a small proportional range for the
coefficients themselves. The maximal extent of that range and
extensions of the half-term remainder law are separate questions:
a coefficient asymptotic alone does not control a normalized
signed-cut remainder.
